\begin{textsample}{POS Dim 3 – ms – Score 11.00 – ms/ms\_inf\_e\_ef\_30.txt}  \label{ex:f3_pos_199}
7.6 Transição para o Ensino Fundamental São os passos que fazem os caminhos.
As reticências são os três primeiros passos do pensamento que continua por conta própria o seu caminho.
Mario Quintana A inserção da \textbf{criança} da Educação \textbf{Infantil} no Ensino Fundamental deve assegurar o seu direito de ser \textbf{criança}, que transita entre o mundo concreto e o imaginário, construindo conceitos de forma \textbf{lúdica} e com liberdade.
A organização de ambientes e práticas educativas para favorecer a aprendizagem e o desenvolvimento deve respeitar o tempo e o modo de aprender e se expressar de cada \textbf{criança}.
Assim, é importante que o trabalho pedagógico ocorra por meio de um planejamento estruturado com objetivos claros e intencionalidade educativa.
Também é necessário considerar que algumas \textbf{crianças} não frequentaram a Educação \textbf{Infantil}, então, o trabalho do professor pode partir de um diagnóstico inicial para saber quais conhecimentos elas já possuem e o que é preciso planejar para atender as suas necessidades.
Nesse sentido, Kramer ( 2007, p.19-20 ) aponta que : > Questões como alfabetizar ou não na educação \textbf{infantil} e ensino fundamental continuam atuais.
Temos \textbf{crianças}, sempre, na educação \textbf{infantil} e no ensino fundamental.
Entender que as pessoas são sujeitos da história e da cultura, além de serem por elas produzidas, e considerar milhões de estudantes brasileiros de 0 a 10 anos como \textbf{crianças} e não só estudantes, implica ver o pedagógico na sua dimensão cultural, como conhecimento, arte e vida, e não só como algo instrucional, que visa a ensinar coisas.
Essa reflexão vale para a educação \textbf{infantil} e o ensino fundamental.
Cabe destacar a importância de se pensar na mudança de etapa como continuidade : tanto a Educação \textbf{Infantil} como os primeiros anos do Ensino Fundamental envolvem \textbf{cuidados}, atenção, conhecimentos, aprendizagens, conquistas e ludicidade, devido à faixa \textbf{etária} dessas \textbf{crianças}.
Essa passagem deve ser tranquila para elas e seus familiares, evitando ansiedade e insegurança.
As \textbf{crianças} constroem conhecimentos, têm direitos que precisam ser assegurados, considerando suas necessidades de aprender, \textbf{brincar}, interagir, experimentar, afinal é imprescindível que o professor favoreça vivências planejadas e significativas.
É importante desenvolver ações articuladas entre as etapas da Educação \textbf{Infantil} e do Ensino Fundamental para que os professores, a equipe pedagógica e a gestão compreendam as especificidades de cada uma, com o propósito de estabelecer estratégias de transição, de forma a respeitar os direitos das \textbf{crianças} e a \textbf{infância}.
Nas Diretrizes Curriculares Nacionais para a Educação \textbf{Infantil} ( DCNEI ), art. 10, é abordada a questão “ de que as \textbf{instituições} devem criar procedimentos de \textbf{acompanhamento} do trabalho pedagógico e para a avaliação do desenvolvimento das \textbf{crianças}, sem objetivo de seleção, promoção ou classificação... ” Assim, os instrumentos avaliativos são \textbf{registros} importantes para acompanhar a aprendizagem e o desenvolvimento das \textbf{crianças}.
Sobre essa vertente, a BNCC dispõe que : > [... ] as informações contidas em relatórios, portfólios ou outros \textbf{registros} que evidenciem os processos vivenciados pelas \textbf{crianças} ao longo de sua trajetória na Educação \textbf{Infantil} podem contribuir para a compreensão da história da vida escolar de cada estudante do Ensino Fundamental.
Conversas ou visitas e troca de materiais entre os professores das escolas de Educação \textbf{Infantil} e de Ensino Fundamental – Anos Iniciais também são importantes para facilitar a inserção das \textbf{crianças} nessa nova etapa da vida escolar ( BRASIL, 2017 ).
Ainda no âmbito das DCNEI, o art. 11 aponta que “ Na transição para o Ensino Fundamental a proposta pedagógica deve prever formas para garantir a continuidade no processo de aprendizagem e desenvolvimento das \textbf{crianças}, respeitando as especificidades \textbf{etárias}, sem antecipação de conteúdos que serão trabalhados no Ensino Fundamental ”.
Portanto, é fundamental que os objetivos da Educação \textbf{Infantil} estejam voltados para as necessidades das \textbf{crianças} dessa faixa \textbf{etária}, sem adiantar nenhuma etapa.
Sobre a inserção das \textbf{crianças} em uma nova etapa da Educação Básica, a BNCC recomenda : > A transição entre essas duas etapas da Educação Básica requer muita atenção para que haja \textbf{equilíbrio} entre as mudanças introduzidas, garantindo integração e continuidade dos processos de aprendizagens das \textbf{crianças}, respeitando suas singularidades e as diferentes relações que elas estabelecem com os conhecimentos, assim como a natureza das mediações de cada etapa.
Torna -se necessário estabelecer estratégias de acolhimento e adaptação tanto para as \textbf{crianças} quanto para os docentes, de modo que a nova etapa se construa com base no que a \textbf{criança} sabe e é capaz de fazer em uma perspectiva de continuidade de seu percurso educativo ( BRASIL, 2017 ).
Para que a recomendação acima se efetive, é importante o respeito às ideias, às explicações pessoais e à forma de pensar das \textbf{crianças} que estão em transição para uma nova etapa da Educação Básica.

% matched lemmas: acompanhamento, brincar, criança, cuidado, equilíbrio, etário, infantil, infância, instituição, lúdico, registro
\end{textsample}
